\documentclass[10pt,a4paper]{amsart}
\usepackage[margin=19mm]{geometry}
\usepackage{amsmath,amssymb,mathtools}
\usepackage[scr=rsfs,cal=euler]{mathalfa}
\usepackage{xfrac}
\usepackage[unicode,hidelinks,bookmarks=false]{hyperref}
\newcommand{\ie}{i.e.\ }
\newcommand{\cf}{cf.\ }
\newcommand{\resp}{respectively\ }
\newcommand{\Subsection}{\subsection}
\providecommand{\nobreakdash}{\nobreak\hbox{-}}
\setlength{\emergencystretch}{2em}
\multlinegap0pt
\usepackage{amssymb}
\usepackage{enumerate}
\usepackage{algorithm}\usepackage{algpseudocode}
\newtheorem{theorem}{Theorem}
\newtheorem{lemma}{Lemma}
\newtheorem{definition}{Definition}
\newtheorem{case}{Case}
\newtheorem{claimA}{Claim}
\title{\Large \bf The absence of monochromatic triangle implies various properly colored spanning trees}
\author{Ruonan Li, Ruhui Lu, Xueli Su, Shenggui Zhang}
\date{Source: arXiv:2403.09082v3 (CC BY 4.0)}
\begin{document}
\maketitle

\noindent\emph{Source and licence.} \Large \bf The absence of monochromatic triangle implies various properly colored spanning trees. Version arXiv:2403.09082v3 (CC BY 4.0), distributed by arXiv under the Creative Commons Attribution 4.0 International licence, http://creativecommons.org/licenses/by/4.0/, as recorded in the arXiv licence metadata for this exact version and shown on the abstract page https://arxiv.org/abs/2403.09082v3. Full licence notice: \texttt{licenses/arxiv-2403.09082-CC-BY-4.0.txt}.\par\smallskip

\noindent\textit{Editorial scope.} Counted unit: the article's theorem Thm:k-spider (source0.tex lines 101--104). The article's complete original proof of it is delivered with it (source0.tex lines 411--458, one \texttt{proof} environment of 48 lines, which the article places away from the statement). No mathematics is written, repaired or re-derived: the statement and the proof are the article's own lines, in the article's own notation, and no retained line is substituted or re-flowed.  Retained uncounted as the source-local context the proof needs: lines 175--177; lines 182--184; lines 218--220; lines 241--243; lines 278--280.  Nothing was omitted from any retained span, so the package carries no omission record.  No retained line contains a citation, so the wrapper carries no bibliography and no bibliography entry.  No retained line contains a figure environment, an image-inclusion command or drawing code, so no image is delivered and the package declares no asset dependency.  Editorial formatting only, in addition: an AMS \texttt{amsart} preamble that restates the article's own declarations and macros used by the retained lines, namely the environments \texttt{theorem}, \texttt{lemma}, \texttt{definition}, \texttt{case}, \texttt{claimA} on separate counters, together with the standard packages those lines need.  The retained environments print in this document as Theorem 1, Lemma 1, Definition 1, Case 1, Claima 1 in order of appearance, whereas the article prints the same items with the numbering of its own full document.  The complete native source of the article as distributed in the arXiv e-print for this exact version is bundled unchanged as \texttt{sources/156057/source0.tex}.
\begin{lemma}\label{Lem:STAR}
    Let $G$ be a mono-$C_3$-free edge-colored $K_n$. Then for each vertex $v\in V(G)$, there is a PC Hamilton path in $G$ starting at $v$.
\end{lemma}

\begin{lemma}[Spanning nice shovel]\label{Lem:SPANING}
     Let $G$ be a mono-$C_3$-free edge-colored $K_n$ with $n\geq 3$. Then $G$ contains a spanning subgraph which is a nice $(n-2)$-shovel.
\end{lemma}

\begin{lemma}\label{Lem:S AND Y}
Let $G$ be a mono-$C_3$-free edge-colored $K_n$. Let $S$ and $Y$ be vertex-disjoint subgraphs of $G$ such that $S=u_1u_2u_3u_1$ is a triangle with $u_2u_3$ being a center edge and $Y$ is a nice shovel or a vertex set of size at most 2. Then $G$ contains a PC path $P=z_1z_2\cdots z_\ell$ such that $V(P)=V(S)\cup V(Y)$, $z_1=u_1$ and $col(z_1z_2)\in col(u_1,\{u_2,u_3\})$.
\end{lemma}

\begin{definition}\label{def:monoC3-free-tournament}
    Given $t\geq 2$, let  $D$ be a multipartite tournament such that each partite set has at most 2 vertices and $N^+\left( x \right) \cap N^+\left( y \right) =\emptyset$ for every pair of distinct vertices $x$ and $y$ from a same partite set. Then we say $D$ is a {\it mono-$C_3$-free tournament}.
\end{definition}

\begin{theorem}\label{Thm:D-spider}
  For $k\ge 2$ and $\ell_2\geq \ell_3 \geq \cdots \geq \ell_k\geq 1$,  let $D$ be a mono-$C_3$-free tournament on $n\geq\sum^k_{i=2}{\ell_i}+2k^2+2k+6$ vertices. If $|N^+_D(v)|\geq 2$  for all $v\in V(D)$.  Then $D$ contains  a  subgraph isomorphic to   $T^*_{\ell_2,\ldots,\ell_k}$.
\end{theorem}

\begin{theorem}\label{Thm:k-spider}
     Given an integer $k\geq 3$, let $G$ be a mono-$C_3$-free edge-colored $K_n$. If $n\geq 6k\cdot g(S_k,C_3)+2k^3+2k^2+8k$,  then
    for every $k$ positive integers $\ell_1\geq \ell_2\geq \cdots\geq \ell_k$ satisfying $\sum_{i=1}^k \ell_i=n-1$, $G$ contains a properly colored spanning spider $T$ with legs of lengths $\ell_1,\ell_2,\ldots,\ell_k$, respectively.
\end{theorem}

\begin{proof}[\textbf{Proof of Theorem \ref{Thm:k-spider}}]
The proof is given by distinguishing the number of vertex-disjoint nice bowties in $G$.
\begin{case}
    There are at least $g(S_k,C_3)+1$ vertex-disjoint nice bowties in $G$.
\end{case}
Let $B_1,\cdots,B_s$ be vertex-disjoint nice bowties in $G$ with $s\geq g(S_k,C_3)+1$ such that $G-\cup^h_{i=1} {V(B_i)}$  contains no nice bowtie.   Pick one center $v_i$ from each bowtie $B_i$ ($1\leq i\leq s$) to form a set  $W$.  Since $G$ is mono-$C_3$-free and $|W|\geq g(S_k,C_3)+1$, there exists a rainbow $k$-star $S_k$ in $G[W]$. Without loss of generality, assume $V(S_k)=\{v_0, v_1,v_2, \ldots, v_k\}$ and $v_0$ is the $k$-degree vertex in $S_k$.

For $1\leq i\leq k$, since $v_i$ is a center vertex of the bowtie $B_i$, there exists a subgraph $X_i\subseteq B_i$ such that $X_i+v_iv_0$ is either a nice 2-shovel or a nice 3-shovel. Denote by $Y_i$  the nice shovel $X_i+v_iv_0$. Then $Y_1,Y_2,\ldots,Y_k$ are disjoint shovels overlapping at the vertex $v_0$. We will obtain a desired PC spider by extending or shrinking these shovels.

Define $I=\{1,2,\ldots,k\}$ and $I^-=\{i\in I:  \ell_i \leq |Y_i|-1\}$.
For each $i\in I^-$,  we can modify $Y_i$ into a leg of length $\ell_i$ by removing  $|Y_i|-1-\ell_i $ vertices.  Let $R$  be the removed vertices. Then $U=(V(G)\setminus \cup^k_{i=1} V(Y_i))\cup R$ is the set of vertices waiting to be embedded. Note that $|U|=\sum_{i\in I\setminus I^-} {(\ell_i-|Y_i|+1)}$. We partition $U$ into $|I\setminus I^-|$ disjoint sets $\{U_i\}_{i\in I\setminus I^-}$  such that $|U_i|=\ell_i-|Y_i|+1$ for each $i\in I\setminus I^-$.
Then by Lemma \ref{Lem:SPANING}, each $G[U_i]$ of order at least $3$ contains a spanning nice shovel $F_i$. Set $F_i=U_i$ when $|U_i|\leq 2$. Apply Lemma \ref{Lem:S AND Y} to the triangle in $Y_i$ and $F_i$ for all $i\in I\setminus I^-$, we get a desired PC spanning tree in $G$.

\begin{case}
   The number of vertex-disjoint nice bowties in $G$ is at most $g(S_k,C_3)$.
\end{case}
In this case we will first find a structure that is almost the spider we want except the first leg, which is a nice 1-shovel. Then we get the final spider by extending this shovel into a leg of length $\ell_1$.
To state the proof, we need more definitions.
Let $S$ be a spider with precisely $k-1$ legs of lengths $\ell_2,\ell_3,\ldots,\ell_k$ respectively, let $Y$ be a triangle and let $O^*_{ \ell_2, \ldots, \ell_k}$ be the graph obtained by identifying a vertex in $Y$ with the center of $S$. We call $O^*_{ \ell_2, \ldots, \ell_k}$ an {\it octopus}. The legs of $O^*_{ \ell_2, \ldots, \ell_k}$ are exactly the legs of $S$. We say $O^*_{ \ell_2, \ldots, \ell_k}$ is {\it nice} if $S$ is a PC spider and $Y\cup P$ is a nice shovel for each leg $P$ of $S$. If $G$ contains a {nice} octopus $O^*_{ \ell_2, \ldots, \ell_k}$, then let $F=G-O^*_{ \ell_2, \ldots, \ell_k}$.
By Lemma \ref{Lem:SPANING},  $F$ contains a spanning nice shovel $F^\prime$ when $|F|\geq 3$. Set $F^\prime=F$ when $|F|\leq 2$. Apply Lemma \ref{Lem:S AND Y} to $F^\prime$ and the triangle in $O^*_{ \ell_2, \ldots, \ell_k}$, we get a desired PC spider. Therefore the following claim holds.
\begin{claimA}\label{Cl:octopus}
	If $G$ contains a nice octopus $O^*_{ \ell_2, \ldots, \ell_k}$, then $G$ contains a PC spider with legs of lengths $\ell_1,\ell_2, \ldots, \ell_k$.
\end{claimA}


Let $B_1,B_2,\ldots,B_s$ be vertex-disjoint nice bowties in $G$ such that there is no nice bowtie in $H:=G\setminus \cup_{i=1}^s V(B_i)$. Then $s\leq g(S_k,C_3)$. Recall that each bowtie has at most $6$ vertices. We have $|H|\geq n-6g(S_k,C_3)$. Note that $\ell_1\geq \frac{n-1}{k}\geq 6g(S_k,C_3)+2k^2+2k+7$. Then
$$|H|=1+\sum_{i=2}^k\ell_i+\ell_1-6g(S_k,C_3)\geq \sum_{i=2}^k\ell_i+2k^2+2k+8.$$

For each vertex $v$ in $H$, we claim that at most one color in $col(v,N_{H}(v))$ appears more than once at $v$. Otherwise by the mono-$C_3$-free condition, $H$  contains a nice short bowtie with $v$ being the center vertex, which contradicts the choice of $H$.

\begin{claimA}\label{Cl:color-once}
For each vertex $v\in V(H)$, at most one color in $col(v,N_{H}(v))$ appears more than once at $v$.
\end{claimA}

There could be a vertex $v$ with no color repeats at $v$. Let $V_1=\{v\in V(H): \Delta^{mon}_H(v)=1\}$ and $V_2=\{v\in V(H): \Delta^{mon}_H(v)\geq 2\}$.

If $|V_1|\geq 3$, then let $v_1,v_2,v_3$ be distinct vertices in $V_1$. By the definition of $V_1$, the cycle $v_1v_2v_3v_1$ must be a rainbow triangle. Let $H^\prime=H-\{v_1,v_2,v_3\}$. Recall that $|H|\geq \sum_{i=2}^k\ell_i+2k^2+2k+8$. We have $|H^\prime|\geq \sum_{i=2}^k\ell_i+2k^2+2k+5$. Let $X_2, X_3,\ldots, X_k$ be disjoint sets in $H^\prime$ such that $|X_i|=\ell_i$ for $2\leq i\leq k$. By Lemma \ref{Lem:STAR}, there is a PC path $P_i$ on $\{v_1\}\cup X_i$ with $v_1$ being the starting vertex for $2\leq i\leq k$. Since $\Delta^{mon}_{H^\prime}(v_1)=1$, the triangle $v_1v_2v_3v_1$ and paths $P_2,P_3,\ldots,P_k$ form a nice octopus $O^*_{ \ell_2, \ldots, \ell_k}$. By Claim \ref{Cl:octopus}, we are home.

If $2\leq \Delta^{mon}_{H}(x)\leq 2k^2+2k+7$ for some vertex $x\in V_2$, then assume $\alpha$ is the color such that $\Delta^{mon}_{H}(x) = |\{u\in V(H): col(ux) = \alpha\}|$.  Let $U_\alpha=\{u\in V(H): col(ux) = \alpha\}$ and $H^\prime=H-U_\alpha$. Then $|H^\prime|\geq \sum_{i=2}^k\ell_i+1$. Let $X_2, X_3,\ldots, X_k$ be disjoint sets in $H^\prime-x$ such that $|X_i|=\ell_i$ for $2\leq i\leq k$. By Lemma \ref{Lem:STAR}, there is a PC path $P_i$ on $\{x\}\cup X_i$ with $v_1$ being the starting vertices for $2\leq i\leq k$. Note that $\Delta^{mon}_{H^\prime}(x)=1$. Choose distinct vertices $a,b\in U_\alpha$, then the triangle $xabx$ and paths $P_2,P_3,\ldots,P_k$ form a nice octopus $O^*_{ \ell_2, \ldots, \ell_k}$. By Claim \ref{Cl:octopus}, we are home.

The remaining case is that $|V_1|\leq 2$ and $\Delta^{mon}_{H}(x)\geq 2k^2+2k+8$ for every $x\in V_2$. Let $H_1=H-V_1$.
Then $|H_1|\geq \sum_{i=2}^k\ell_i+2k^2+2k+6$ and $\Delta^{mon}_{H_1}(u)\geq 2k^2+2k+6$ for every vertex $u\in V(H_1)$. Let $f(u)$ be the unique color appearing more than once in $col(u,N_{H_1}(u))$. We claim that for any two vertices $u, v\in V(H_1)$,  there holds $col(u)=f(u)$ or $f(v)$. Otherwise, $H$ contains a nice long bowtie with $u$ and $v$ being the center vertices, which contradicts the choice of $H$.
\begin{claimA}\label{Cl:degenerate}
For each pair of vertices $u,v\in V(H)$, we have $col(uv)=f(u)$ or $f(v)$.
\end{claimA}
Now construct an auxiliary digraph $D$ satisfying $V(D)=V(H_1)$ and $A(D)=\{uv: col(uv)=f(u) \ \text{and}\  col(uv)\ne f(v) \}.$ Since $H_1$ is mono-$C_3$-free, according to Claim \ref{Cl:degenerate} and Definition \ref{def:monoC3-free-tournament}, the directed graph $D$ is a mono-$C_3$-free tournament with $N^+_{D}(u)\geq 2k^2+2k+6>2$ for every vertex $u\in V(H_1)$. Recall that $|H_1|\geq \sum_{i=2}^k\ell_i+2k^2+2k+6$.
By Theorem \ref{Thm:D-spider}, $D$ contains a subgraph $ T^*_{\ell_2, \cdots, \ell_k}$, which is a nice octopus $O^*_{ \ell_2,\ldots,\ell_k}$ in $G$. Again we obtain the desired PC spanning tree by Claim \ref{Cl:octopus}. The proof is complete.
\end{proof}

\end{document}
